\section{Reason for Choosing the Topic}

In recent years, Legal AI has emerged as a transformative area to assist legal professionals in parsing and applying massive statutory databases. While proprietary Large Language Models (LLMs) excel at zero-shot statutory interpretation and generating rationales, deploying them in institutional legal workflows is constrained by critical challenges:
\begin{itemize}
    \item \textbf{Data Privacy:} Transmitting sensitive legal queries containing personal or corporate data to external APIs of commercial providers violates localized data sovereignty and regulations like the Cybersecurity Law of Vietnam.
    \item \textbf{Operating Costs:} Recurring API token charges for processing long statutory contexts are financially unsustainable for public courts and local firms.
    \item \textbf{Availability:} API-dependent applications are subject to network latency and downtime, which are unacceptable for high-throughput institutional processing.
\end{itemize}

Deploying lightweight Small Language Models (SLMs) with fewer than 3 billion parameters on-premise offers a secure, cost-effective alternative. However, SLMs suffer from a severe "reasoning gap"---they exhibit weak multi-step deduction, frequently hallucinate legal clauses, and suffer from logic loops. In legal tasks, a single logical error can render an entire verdict invalid \parencite{li2026legaldrilldiagnosisdrivensynthesislegal, stacey2024distillingrobustnessnaturallanguage}.

To bridge this gap, knowledge distillation (KD) transfers domain knowledge from a reinforcement-learning-aligned teacher to a student model \parencite{hinton2015distillingknowledgeneuralnetwork}. Yet, standard offline distillation suffers from exposure bias during sequential text generation and fails to address the student's inner calibration blind spots. 

To address these limitations, this study proposes a hybrid framework that integrates offline logit distillation, on-policy trajectory refinement to counteract exposure bias, and diagnosis-driven preference optimization to correct systematic reasoning fallacies. By tailoring this secure pipeline to Vietnamese legal tasks (Natural Language Inference, Multiple-Choice QA, and Syllogism), we demonstrate the feasibility of local, high-performing legal assistants, which serves as the core motivation for this research topic.

